\documentclass[10pt,a4paper]{amsart}
\usepackage{amsmath,amssymb,amsthm,graphicx}
\usepackage{verbatim}
\newenvironment{reference}{\comment}{\endcomment}
%\newenvironment{reference}{}{}
\newenvironment{slogan}{\comment}{\endcomment}
\newenvironment{history}{\comment}{\endcomment}

% For commutative diagrams we use Xy-pic
\usepackage[all]{xy}

% We use 2cell for 2-commutative diagrams.
\xyoption{2cell}
% Theorem environments.
%
\theoremstyle{plain}
\newtheorem{theorem}[subsection]{Theorem}
\newtheorem{proposition}[subsection]{Proposition}
\newtheorem{lemma}[subsection]{Lemma}

\theoremstyle{definition}
\newtheorem{definition}[subsection]{Definition}
\newtheorem{example}[subsection]{Example}
\newtheorem{exercise}[subsection]{Exercise}
\newtheorem{situation}[subsection]{Situation}

\theoremstyle{remark}
\newtheorem{remark}[subsection]{Remark}
\newtheorem{remarks}[subsection]{Remarks}

\numberwithin{equation}{subsection}

% Macros
%
\def\lim{\mathop{\mathrm{lim}}\nolimits}
\def\colim{\mathop{\mathrm{colim}}\nolimits}
\def\Spec{\mathop{\mathrm{Spec}}}
\def\Hom{\mathop{\mathrm{Hom}}\nolimits}
\def\Ext{\mathop{\mathrm{Ext}}\nolimits}
\def\SheafHom{\mathop{\mathcal{H}\!\mathit{om}}\nolimits}
\def\SheafExt{\mathop{\mathcal{E}\!\mathit{xt}}\nolimits}
\def\Sch{\mathit{Sch}}
\def\Mor{\mathop{\mathrm{Mor}}\nolimits}
\def\Ob{\mathop{\mathrm{Ob}}\nolimits}
\def\Sh{\mathop{\mathit{Sh}}\nolimits}
\def\NL{\mathop{N\!L}\nolimits}
\def\CH{\mathop{\mathrm{CH}}\nolimits}
\def\proetale{{pro\text{-}\acute{e}tale}}
\def\etale{{\acute{e}tale}}
\def\QCoh{\mathit{QCoh}}
\def\Ker{\mathop{\mathrm{Ker}}}
\def\Im{\mathop{\mathrm{Im}}}
\def\Coker{\mathop{\mathrm{Coker}}}
\def\Coim{\mathop{\mathrm{Coim}}}

% Boxtimes
%
\DeclareMathSymbol{\boxtimes}{\mathbin}{AMSa}{"02}

%
% Macros for moduli stacks/spaces
%
\def\QCohstack{\mathcal{QC}\!\mathit{oh}}
\def\Cohstack{\mathcal{C}\!\mathit{oh}}
\def\Spacesstack{\mathcal{S}\!\mathit{paces}}
\def\Quotfunctor{\mathrm{Quot}}
\def\Hilbfunctor{\mathrm{Hilb}}
\def\Curvesstack{\mathcal{C}\!\mathit{urves}}
\def\Polarizedstack{\mathcal{P}\!\mathit{olarized}}
\def\Complexesstack{\mathcal{C}\!\mathit{omplexes}}
% \Pic is the operator that assigns to X its picard group, usage \Pic(X)
% \Picardstack_{X/B} denotes the Picard stack of X over B
% \Picardfunctor_{X/B} denotes the Picard functor of X over B
\def\Pic{\mathop{\mathrm{Pic}}\nolimits}
\def\Picardstack{\mathcal{P}\!\mathit{ic}}
\def\Picardfunctor{\mathrm{Pic}}
\def\Deformationcategory{\mathcal{D}\!\mathit{ef}}
\usepackage{cleveref}
\crefname{originalreference}{original reference}{original references}
\Crefname{originalreference}{Original reference}{Original references}
\makeatletter
\newlabel{situation-shriek}{{0F42}{source}{}{sourcepointer-1}{}}
\newlabel{situation-shriek@cref}{{[originalreference][0][]0F42}{[1][0][]source}{}{}{}}
\makeatother
\crefname{originaltheorem}{original theorem}{original theorems}
\Crefname{originaltheorem}{Original theorem}{Original theorems}
\makeatletter
\newlabel{flat-theorem-nagata}{{0F41}{source}{}{sourcepointer-2}{}}
\newlabel{flat-theorem-nagata@cref}{{[originaltheorem][0][]0F41}{[1][0][]source}{}{}{}}
\makeatother
\crefname{originallemma}{original lemma}{original lemmas}
\Crefname{originallemma}{Original lemma}{Original lemmas}
\makeatletter
\newlabel{flat-lemma-compactifyable}{{0A9Z}{source}{}{sourcepointer-3}{}}
\newlabel{flat-lemma-compactifyable@cref}{{[originallemma][0][]0A9Z}{[1][0][]source}{}{}{}}
\makeatother
\crefname{originallemma}{original lemma}{original lemmas}
\Crefname{originallemma}{Original lemma}{Original lemmas}
\makeatletter
\newlabel{lemma-twisted-inverse-image}{{0A9E}{source}{}{sourcepointer-4}{}}
\newlabel{lemma-twisted-inverse-image@cref}{{[originallemma][0][]0A9E}{[1][0][]source}{}{}{}}
\makeatother
\crefname{originalequation}{original equation}{original equations}
\Crefname{originalequation}{Original equation}{Original equations}
\makeatletter
\newlabel{equation-sheafy}{{0A9L}{source}{}{sourcepointer-5}{}}
\newlabel{equation-sheafy@cref}{{[originalequation][0][]0A9L}{[1][0][]source}{}{}{}}
\makeatother
\crefname{originallemma}{original lemma}{original lemmas}
\Crefname{originallemma}{Original lemma}{Original lemmas}
\makeatletter
\newlabel{lemma-proper-noetherian}{{0A9P}{source}{}{sourcepointer-6}{}}
\newlabel{lemma-proper-noetherian@cref}{{[originallemma][0][]0A9P}{[1][0][]source}{}{}{}}
\makeatother
\crefname{originallemma}{original lemma}{original lemmas}
\Crefname{originallemma}{Original lemma}{Original lemmas}
\makeatletter
\newlabel{flat-lemma-compactifications-cofiltered}{{0ATU}{source}{}{sourcepointer-7}{}}
\newlabel{flat-lemma-compactifications-cofiltered@cref}{{[originallemma][0][]0ATU}{[1][0][]source}{}{}{}}
\makeatother
\crefname{originallemma}{original lemma}{original lemmas}
\Crefname{originallemma}{Original lemma}{Original lemmas}
\makeatletter
\newlabel{flat-lemma-right-multiplicative-system}{{0ATV}{source}{}{sourcepointer-8}{}}
\newlabel{flat-lemma-right-multiplicative-system@cref}{{[originallemma][0][]0ATV}{[1][0][]source}{}{}{}}
\makeatother
\crefname{originallemma}{original lemma}{original lemmas}
\Crefname{originallemma}{Original lemma}{Original lemmas}
\makeatletter
\newlabel{categories-lemma-morphisms-right-localization}{{04VI}{source}{}{sourcepointer-9}{}}
\newlabel{categories-lemma-morphisms-right-localization@cref}{{[originallemma][0][]04VI}{[1][0][]source}{}{}{}}
\makeatother
\crefname{originallemma}{original lemma}{original lemmas}
\Crefname{originallemma}{Original lemma}{Original lemmas}
\makeatletter
\newlabel{categories-lemma-equality-morphisms-right-localization}{{04VJ}{source}{}{sourcepointer-10}{}}
\newlabel{categories-lemma-equality-morphisms-right-localization@cref}{{[originallemma][0][]04VJ}{[1][0][]source}{}{}{}}
\makeatother
\crefname{originallemma}{original lemma}{original lemmas}
\Crefname{originallemma}{Original lemma}{Original lemmas}
\makeatletter
\newlabel{lemma-shriek-well-defined}{{0AA0}{source}{}{sourcepointer-11}{}}
\newlabel{lemma-shriek-well-defined@cref}{{[originallemma][0][]0AA0}{[1][0][]source}{}{}{}}
\makeatother
\crefname{originallemma}{original lemma}{original lemmas}
\Crefname{originallemma}{Original lemma}{Original lemmas}
\makeatletter
\newlabel{lemma-compose-base-change-maps}{{0ATQ}{source}{}{sourcepointer-12}{}}
\newlabel{lemma-compose-base-change-maps@cref}{{[originallemma][0][]0ATQ}{[1][0][]source}{}{}{}}
\makeatother
\crefname{originallemma}{original lemma}{original lemmas}
\Crefname{originallemma}{Original lemma}{Original lemmas}
\makeatletter
\newlabel{lemma-compose-base-change-maps-horizontal}{{0ATR}{source}{}{sourcepointer-13}{}}
\newlabel{lemma-compose-base-change-maps-horizontal@cref}{{[originallemma][0][]0ATR}{[1][0][]source}{}{}{}}
\makeatother
\crefname{originalreference}{original reference}{original references}
\Crefname{originalreference}{Original reference}{Original references}
\makeatletter
\newlabel{remark-going-around}{{0ATS}{source}{}{sourcepointer-14}{}}
\newlabel{remark-going-around@cref}{{[originalreference][0][]0ATS}{[1][0][]source}{}{}{}}
\makeatother
\title{The upper shriek functors of composable morphisms compose}
\author{The Stacks Project Authors}
\date{stacks-project:duality.tex:local-snapshot-20261012; GNU FDL 1.2 or later}
\begin{document}
\maketitle
\noindent\textit{Editorial note.} All delivered mathematical text is the author's own, copied from the pinned Stacks Project chapter file duality.tex; no statement or proof step was rewritten, reordered or removed.\par\medskip
\noindent\textit{Editorial note.} Preamble provenance: the pinned chapter file opens with an input of the shared preamble and carries no preamble text of its own, so the project's single Stacks-wide preamble staged at sources/stacks-duality/preamble.tex is excerpted verbatim under the recipe fixed by LANE\_BRIEF 22.1, as three spans, lines 8-13 (verbatim package and the reference/slogan/history environments), lines 14-18 (Xy-pic setup) and lines 154-217 (the native theorem declarations and mathematical macros). Nothing was added, removed or redefined; the wrapper class and the amsmath/amssymb/amsthm/graphicx base packages are supplied by the assembler.\par\medskip
\noindent\textit{Editorial note.} The counted statement and proof contain no cross reference, no citation and no included file, so the delivered body is self contained.\par\medskip
\noindent\textit{Editorial note.} Canonical location of the counted proposition: Stacks tag 0ATX (sources/stacks-duality/pinned-original-source.tex line 4161).\par\medskip
\noindent\textit{Allowed standard prerequisites (editorial).}\begin{itemize}
\item The upper shriek functor as the right adjoint of the twisted inverse image, and its definition through a compactification. Same chapter, Lemma labelled lemma-twisted-inverse-image and Situation labelled situation-shriek, cited at proof lines 4167-4184. Supplies the functors being composed and the compactification description of each.
\item Nagata compactification and the existence of compactifications. More on Flatness, Theorem labelled flat-theorem-nagata and Lemma labelled flat-lemma-compactifyable, cited at proof lines 4169-4171. Produces the compactifications used in the construction.
\item The base change map between the right adjoints for a cartesian square, and its independence of choices. Same chapter, equation labelled equation-sheafy, cited at proof lines 4188-4192, and Lemma labelled lemma-proper-noetherian, cited at proof lines 4198-4200. Gives the map gamma and shows that it is an isomorphism on D\_QCoh\textasciicircum{}+.
\item Cofilteredness of compactifications and composition rules for base change maps. More on Flatness, Lemmas labelled flat-lemma-compactifications-cofiltered and flat-lemma-right-multiplicative-system, cited at proof lines 4218-4226; Same chapter, Lemmas labelled lemma-compose-base-change-maps and lemma-compose-base-change-maps-horizontal, cited at proof lines 4304-4308, with Remark labelled remark-going-around. Reduces independence of choices to one diagram and lets the five gamma maps be reordered.
\item Localisation of the category of morphisms with respect to right multiplicative systems. Categories, Lemmas labelled categories-lemma-morphisms-right-localization and categories-lemma-equality-morphisms-right-localization, cited at proof lines 4226-4230. Allows the compatible choice of the compactification of X.
\end{itemize}\medskip
\ifdefined\hypertarget\hypertarget{sourcepointer-1}{}\fi
\noindent\textit{Source reference.} Original reference 0F42: sources/stacks-duality/pinned-original-source.tex, source lines 3953--3953. Same-chapter pointer into the author environment situation at source line 3953; invoked at proof line None of the counted proof. It is retained only as a navigation pointer; every load-bearing step is inside the delivered span.\par\medskip
\ifdefined\hypertarget\hypertarget{sourcepointer-2}{}\fi
\noindent\textit{Source reference.} Original theorem 0F41: sources/stacks-chow/tags.txt, source lines 18886--18886. Cross-chapter standard prerequisite located in the project tag list staged in the local snapshot; invoked at proof line 4170 of the counted proof. It is retained only as a navigation pointer; every load-bearing step is inside the delivered span.\par\medskip
\ifdefined\hypertarget\hypertarget{sourcepointer-3}{}\fi
\noindent\textit{Source reference.} Original lemma 0A9Z: sources/stacks-chow/tags.txt, source lines 12967--12967. Cross-chapter standard prerequisite located in the project tag list staged in the local snapshot; invoked at proof line 4171 of the counted proof. It is retained only as a navigation pointer; every load-bearing step is inside the delivered span.\par\medskip
\ifdefined\hypertarget\hypertarget{sourcepointer-4}{}\fi
\noindent\textit{Source reference.} Original lemma 0A9E: sources/stacks-duality/pinned-original-source.tex, source lines 653--653. Same-chapter pointer into the author environment lemma at source line 653; invoked at proof line 4173 of the counted proof. It is retained only as a navigation pointer; every load-bearing step is inside the delivered span.\par\medskip
\ifdefined\hypertarget\hypertarget{sourcepointer-5}{}\fi
\noindent\textit{Source reference.} Original equation 0A9L: sources/stacks-duality/pinned-original-source.tex, source lines 988--988. Same-chapter pointer into the author environment equation at source line 988; invoked at proof line 4199 of the counted proof. It is retained only as a navigation pointer; every load-bearing step is inside the delivered span.\par\medskip
\ifdefined\hypertarget\hypertarget{sourcepointer-6}{}\fi
\noindent\textit{Source reference.} Original lemma 0A9P: sources/stacks-duality/pinned-original-source.tex, source lines 1106--1106. Same-chapter pointer into the author environment lemma at source line 1106; invoked at proof line 4209 of the counted proof. It is retained only as a navigation pointer; every load-bearing step is inside the delivered span.\par\medskip
\ifdefined\hypertarget\hypertarget{sourcepointer-7}{}\fi
\noindent\textit{Source reference.} Original lemma 0ATU: sources/stacks-chow/tags.txt, source lines 13627--13627. Cross-chapter standard prerequisite located in the project tag list staged in the local snapshot; invoked at proof line 4234 of the counted proof. It is retained only as a navigation pointer; every load-bearing step is inside the delivered span.\par\medskip
\ifdefined\hypertarget\hypertarget{sourcepointer-8}{}\fi
\noindent\textit{Source reference.} Original lemma 0ATV: sources/stacks-chow/tags.txt, source lines 13628--13628. Cross-chapter standard prerequisite located in the project tag list staged in the local snapshot; invoked at proof line 4235 of the counted proof. It is retained only as a navigation pointer; every load-bearing step is inside the delivered span.\par\medskip
\ifdefined\hypertarget\hypertarget{sourcepointer-9}{}\fi
\noindent\textit{Source reference.} Original lemma 04VI: sources/stacks-chow/tags.txt, source lines 6326--6326. Cross-chapter standard prerequisite located in the project tag list staged in the local snapshot; invoked at proof line 4236 of the counted proof. It is retained only as a navigation pointer; every load-bearing step is inside the delivered span.\par\medskip
\ifdefined\hypertarget\hypertarget{sourcepointer-10}{}\fi
\noindent\textit{Source reference.} Original lemma 04VJ: sources/stacks-chow/tags.txt, source lines 6327--6327. Cross-chapter standard prerequisite located in the project tag list staged in the local snapshot; invoked at proof line 4237 of the counted proof. It is retained only as a navigation pointer; every load-bearing step is inside the delivered span.\par\medskip
\ifdefined\hypertarget\hypertarget{sourcepointer-11}{}\fi
\noindent\textit{Source reference.} Original lemma 0AA0: sources/stacks-duality/pinned-original-source.tex, source lines 3986--3986. Same-chapter pointer into the author environment lemma at source line 3986; invoked at proof line 4319 of the counted proof. It is retained only as a navigation pointer; every load-bearing step is inside the delivered span.\par\medskip
\ifdefined\hypertarget\hypertarget{sourcepointer-12}{}\fi
\noindent\textit{Source reference.} Original lemma 0ATQ: sources/stacks-duality/pinned-original-source.tex, source lines 1238--1238. Same-chapter pointer into the author environment lemma at source line 1238; invoked at proof line 4336 of the counted proof. It is retained only as a navigation pointer; every load-bearing step is inside the delivered span.\par\medskip
\ifdefined\hypertarget\hypertarget{sourcepointer-13}{}\fi
\noindent\textit{Source reference.} Original lemma 0ATR: sources/stacks-duality/pinned-original-source.tex, source lines 1375--1375. Same-chapter pointer into the author environment lemma at source line 1375; invoked at proof line 4337 of the counted proof. It is retained only as a navigation pointer; every load-bearing step is inside the delivered span.\par\medskip
\ifdefined\hypertarget\hypertarget{sourcepointer-14}{}\fi
\noindent\textit{Source reference.} Original reference 0ATS: sources/stacks-duality/pinned-original-source.tex, source lines 1471--1471. Same-chapter pointer into the author environment remark at source line 1471; invoked at proof line 4339 of the counted proof. It is retained only as a navigation pointer; every load-bearing step is inside the delivered span.\par\medskip
\begin{lemma}
\label{lemma-upper-shriek-composition}
In Situation \ref{situation-shriek} let $f : X \to Y$ and $g : Y \to Z$
be composable morphisms of $\textit{FTS}_S$. Then there is a canonical
isomorphism $(g \circ f)^! \to f^! \circ g^!$.
\end{lemma}

\begin{proof}
Choose a compactification $i : Y \to \overline{Y}$ of $Y$ over $Z$.
Choose a compactification $X \to \overline{X}$ of $X$ over
$\overline{Y}$. This uses More on Flatness, Theorem \ref{flat-theorem-nagata}
and Lemma \ref{flat-lemma-compactifyable} twice.
Let $\overline{a}$ be the
right adjoint of Lemma \ref{lemma-twisted-inverse-image} for
$\overline{X} \to \overline{Y}$ and let $\overline{b}$
be the
right adjoint of Lemma \ref{lemma-twisted-inverse-image} for
$\overline{Y} \to Z$.
Then $\overline{a} \circ \overline{b}$ is the
right adjoint of Lemma \ref{lemma-twisted-inverse-image} for
the composition $\overline{X} \to Z$.
Hence $g^! = i^* \circ \overline{b}$ and
$(g \circ f)^! = (X \to \overline{X})^* \circ \overline{a} \circ \overline{b}$.
Let $U$ be the inverse image of $Y$ in $\overline{X}$
so that we get the commutative diagram
$$
\xymatrix{
X \ar[r]_j \ar[d] & U \ar[dl] \ar[r]_{j'} & \overline{X} \ar[dl] \\
Y \ar[r]_i \ar[d] & \overline{Y} \ar[dl] \\
Z
}
$$
Let $\overline{a}'$ be the
right adjoint of Lemma \ref{lemma-twisted-inverse-image} for
$U \to Y$.
Then $f^! = j^* \circ \overline{a}'$. We obtain
$$
\gamma : (j')^* \circ \overline{a} \to \overline{a}' \circ i^*
$$
by (\ref{equation-sheafy}) and we can use it to define
$$
(g \circ f)^! =
(j' \circ j)^* \circ \overline{a} \circ \overline{b} =
j^* \circ (j')^* \circ \overline{a} \circ \overline{b}
\to
j^* \circ \overline{a}' \circ i^* \circ \overline{b} =
f^! \circ g^!
$$
which is an isomorphism on objects of $D_\QCoh^+(\mathcal{O}_Z)$ by
Lemma \ref{lemma-proper-noetherian}. To finish the proof we show that
this isomorphism is independent of choices made.

\medskip\noindent
Suppose we have two diagrams
$$
\vcenter{
\xymatrix{
X \ar[r]_{j_1} \ar[d] & U_1 \ar[dl] \ar[r]_{j'_1} & \overline{X}_1 \ar[dl] \\
Y \ar[r]_{i_1} \ar[d] & \overline{Y}_1 \ar[dl] \\
Z
}
}
\quad\text{and}\quad
\vcenter{
\xymatrix{
X \ar[r]_{j_2} \ar[d] & U_2 \ar[dl] \ar[r]_{j'_2} & \overline{X}_2 \ar[dl] \\
Y \ar[r]_{i_2} \ar[d] & \overline{Y}_2 \ar[dl] \\
Z
}
}
$$
We can first choose a compactification $i : Y \to \overline{Y}$
with dense image of $Y$ over $Z$ which dominates both
$\overline{Y}_1$ and $\overline{Y}_2$,
see More on Flatness, Lemma \ref{flat-lemma-compactifications-cofiltered}.
By More on Flatness, Lemma \ref{flat-lemma-right-multiplicative-system} and
Categories, Lemmas \ref{categories-lemma-morphisms-right-localization} and
\ref{categories-lemma-equality-morphisms-right-localization}
we can choose a compactification $X \to \overline{X}$ with dense image of
$X$ over $\overline{Y}$ with morphisms $\overline{X} \to \overline{X}_1$
and $\overline{X} \to \overline{X}_2$ and such that the composition
$\overline{X} \to \overline{Y} \to \overline{Y}_1$ is equal to
the composition $\overline{X} \to \overline{X}_1 \to \overline{Y}_1$
and such that the composition
$\overline{X} \to \overline{Y} \to \overline{Y}_2$ is equal to
the composition $\overline{X} \to \overline{X}_2 \to \overline{Y}_2$.
Thus we see that it suffices to compare the maps
determined by our diagrams when we have a commutative diagram
as follows
$$
\xymatrix{
X \ar[rr]_{j_1} \ar@{=}[d] & &
U_1 \ar[d] \ar[ddll] \ar[rr]_{j'_1} & &
\overline{X}_1 \ar[d] \ar[ddll] \\
X \ar'[r][rr]^-{j_2} \ar[d] & &
U_2 \ar'[dl][ddll] \ar'[r][rr]^-{j'_2} & &
\overline{X}_2 \ar[ddll] \\
Y \ar[rr]^{i_1} \ar@{=}[d] & & \overline{Y}_1 \ar[d] \\
Y \ar[rr]^{i_2} \ar[d] & & \overline{Y}_2 \ar[dll] \\
Z
}
$$
and moreover the compactifications $X \to \overline{X}_1$ and
$Y \to \overline{Y}_2$ have dense image.
We use $\overline{a}_i$, $\overline{a}'_i$, $\overline{c}$, and
$\overline{c}'$ for the
right adjoint of Lemma \ref{lemma-twisted-inverse-image} for
$\overline{X}_i \to \overline{Y}_i$, $U_i \to Y$,
$\overline{X}_1 \to \overline{X}_2$, and $U_1 \to U_2$.
Each of the squares
$$
\xymatrix{
X \ar[r] \ar[d] \ar@{}[dr]|A & U_1 \ar[d] \\
X \ar[r] & U_2
}
\quad
\xymatrix{
U_2 \ar[r] \ar[d] \ar@{}[dr]|B & \overline{X}_2 \ar[d] \\
Y \ar[r] & \overline{Y}_2
}
\quad
\xymatrix{
U_1 \ar[r] \ar[d] \ar@{}[dr]|C & \overline{X}_1 \ar[d] \\
Y \ar[r] & \overline{Y}_1
}
\quad
\xymatrix{
Y \ar[r] \ar[d] \ar@{}[dr]|D & \overline{Y}_1 \ar[d] \\
Y \ar[r] & \overline{Y}_2
}
\quad
\xymatrix{
X \ar[r] \ar[d] \ar@{}[dr]|E & \overline{X}_1 \ar[d] \\
X \ar[r] & \overline{X}_2
}
$$
is cartesian (see
More on Flatness, Lemma \ref{flat-lemma-compactifications-cofiltered} part (c)
for A, D, E and recall that $U_i$ is the inverse image of $Y$
by $\overline{X}_i \to \overline{Y}_i$ for B, C) and hence
gives rise to a base change map (\ref{equation-sheafy}) as follows
$$
\begin{matrix}
\gamma_A : j_1^* \circ \overline{c}' \to j_2^* &
\gamma_B : (j_2')^* \circ \overline{a}_2 \to \overline{a}'_2 \circ i_2^* &
\gamma_C : (j_1')^* \circ \overline{a}_1 \to \overline{a}'_1 \circ i_1^* \\
\gamma_D : i_1^* \circ \overline{d} \to i_2^* &
\gamma_E : (j'_1 \circ j_1)^* \circ \overline{c} \to (j'_2 \circ j_2)^*
\end{matrix}
$$
Denote $f_1^! = j_1^* \circ \overline{a}'_1$,
$f_2^! = j_2^* \circ \overline{a}'_2$,
$g_1^! = i_1^* \circ \overline{b}_1$,
$g_2^! = i_2^* \circ \overline{b}_2$,
$(g \circ f)_1^! =
(j_1' \circ j_1)^* \circ \overline{a}_1 \circ \overline{b}_1$, and
$(g \circ f)^!_2 =
(j_2' \circ j_2)^* \circ \overline{a}_2 \circ \overline{b}_2$.
The construction given in the first paragraph of the proof
and in Lemma \ref{lemma-shriek-well-defined} uses
\begin{enumerate}
\item $\gamma_C$ for the map $(g \circ f)^!_1 \to f_1^! \circ g_1^!$,
\item $\gamma_B$ for the map $(g \circ f)^!_2 \to f_2^! \circ g_2^!$,
\item $\gamma_A$ for the map $f_1^! \to f_2^!$,
\item $\gamma_D$ for the map $g_1^! \to g_2^!$, and
\item $\gamma_E$ for the map $(g \circ f)^!_1 \to (g \circ f)^!_2$.
\end{enumerate}
We have to show that the diagram
$$
\xymatrix{
(g \circ f)^!_1 \ar[r]_{\gamma_E} \ar[d]_{\gamma_C} &
(g \circ f)^!_2 \ar[d]_{\gamma_B} \\
f_1^! \circ g_1^! \ar[r]^{\gamma_A \circ \gamma_D} & f_2^! \circ g_2^!
}
$$
is commutative. We will use
Lemmas \ref{lemma-compose-base-change-maps} and
\ref{lemma-compose-base-change-maps-horizontal}
and with (abuse of) notation as in
Remark \ref{remark-going-around} (in particular
dropping $\star$ products with identity transformations
from the notation).
We can write $\gamma_E = \gamma_A \circ \gamma_F$ where
$$
\xymatrix{
U_1 \ar[r] \ar[d] \ar@{}[rd]|F & \overline{X}_1 \ar[d] \\
U_2 \ar[r] & \overline{X}_2
}
$$
Thus we see that
$$
\gamma_B \circ \gamma_E = \gamma_B \circ \gamma_A  \circ \gamma_F
= \gamma_A \circ \gamma_B \circ \gamma_F
$$
the last equality because the two squares $A$ and $B$ only
intersect in one point (similar to the last argument in
Remark \ref{remark-going-around}). Thus it suffices to prove that
$\gamma_D \circ \gamma_C = \gamma_B \circ \gamma_F$.
Since both of these are equal to the map (\ref{equation-sheafy})
for the square
$$
\xymatrix{
U_1 \ar[r] \ar[d] & \overline{X}_1 \ar[d] \\
Y \ar[r] & \overline{Y}_2
}
$$
we conclude.
\end{proof}

\end{document}
